\documentclass[aspectratio=169, 14pt]{beamer}
\usepackage{beamer_preamble}

\begin{document}

% ==== Title ===============================================================
\titleslide{Lesson 5-5 \textperiodcentered\ Period 1 (Quick)}{Function Operations + Composition Order DOK-3 (Standalone)}

% ==== Do Now ==============================================================
\begin{frame}{Do Now --- Profit as a function operation}{Algebra 2 \textperiodcentered\ Lesson 5-5 \textperiodcentered\ Period 1}
\phasebadges{1}{5}{bridge to function ops}
\vspace{0.2cm}
\begin{projcard}{Model \& Discuss --- Grooming USA}{slidegold}
Revenue: \$25 per pet. \ Cost: \$15 per pet \(+\) \$1{,}000/month.

\vspace{0.2cm}
\begin{enumerate}
  \setlength{\itemsep}{0.1cm}
  \item Define \(r(x)\) and \(c(x)\).
  \item Did they profit on 95 pets? On 110?
  \item How could you write ONE function for profit?
\end{enumerate}

\vspace{0.1cm}
\small\textcolor{slidegray}{Think alone 3 min --- write \(P(x) = r(x) - c(x)\) in your notes.}
\end{projcard}
\end{frame}

% ==== Launch ==============================================================
\begin{frame}{Launch --- Function Operations Rules (Ex 1 + Ex 4)}{Algebra 2 \textperiodcentered\ Lesson 5-5 \textperiodcentered\ Period 1}
\phasebadges{2}{12}{teacher models}
\vspace{0.15cm}
\begin{projcard}{Function Operations Rules --- keep on your page}{slidegreen}
\small
\begin{enumerate}
  \setlength{\itemsep}{0pt}
  \item \((f\pm g)(x) = f(x)\pm g(x);\ \ (fg)(x) = f(x)\cdot g(x);\ \ (f/g)(x) = f(x)/g(x),\ g(x)\neq 0\).
  \item Domain of combined function \(=\) intersection of domains, minus new excluded values.
  \item Composition: \((f\circ g)(x) = f(g(x))\) --- apply \(g\) FIRST, then \(f\).
  \item \(f\circ g \neq g\circ f\) in general. \textbf{ORDER MATTERS.}
\end{enumerate}
\end{projcard}

\vspace{-0.1cm}
\begin{projcard}{30-sec partner check}{slidegreen}
Test \(f(g(2))\) vs \(g(f(2))\) with \(f(x)=2x-1\) and \(g(x)=3x\).\\
Are the results the same?
\end{projcard}
\end{frame}

% ==== Explore =============================================================
\begin{frame}{Explore --- TPS + DOK-3 Capstone}{Algebra 2 \textperiodcentered\ Lesson 5-5 \textperiodcentered\ Period 1}
\phasebadges{2-3}{33}{student work}
\small\textcolor{slidegray}{8 items in order. Plan \(\sim\)12 min for Practice \#32 (Music Store). 45-min Wed F: skip \#34 + \#28.}

\vspace{0.1cm}
\begin{projcard}{Try It 1 \textperiodcentered\ Add/subtract functions}{slideblue}
\(f(x)=2x^2+7x-1\), \(g(x)=3-2x\). Find \(f+g\) and \(f-g\).
\end{projcard}
\vspace{-0.12cm}
\begin{projcard}{Try It 4 \textperiodcentered\ Compose functions}{slideblue}
\(f(x)=2x-1\), \(g(x)=3x\). Find \(f(g(2))\) and \(f(g(x))\).
\end{projcard}
\vspace{-0.12cm}
\begin{projcard}{\#20 / \#21 \textperiodcentered\ Operations; \quad \#24 \textperiodcentered\ Domain; \quad \#28 / \#34 \textperiodcentered\ Compose + mixed}{slideblue}
\small See packet. Work in order. Domain intersections matter for \#24.
\end{projcard}
\vspace{-0.12cm}
\begin{projcard}{Practice \#32 \(\bigstar\) DOK-3 ORDER-MATTERS --- Music Store Double Discount}{slideaccent}
\small Two promotions: \(\$5\) off OR \(15\%\) off. Apply to \(\$90\) in BOTH orders. Which order saves more? Write the justification (Part C).
\end{projcard}
\end{frame}

% ==== Share / Summary =====================================================
\begin{frame}{Share / Summary --- Order Matters}{Algebra 2 \textperiodcentered\ Lesson 5-5 \textperiodcentered\ Period 1}
\phasebadges{2}{5}{DOK-3 reveal}
\vspace{0.2cm}
\begin{projcard}{Sentence frame \textperiodcentered\ Practice \#32 Music Store}{slideblue}
``Applying \underline{\hspace{1.5cm}} first gives \underline{\hspace{1.5cm}},
then \underline{\hspace{1.5cm}} gives \underline{\hspace{1.5cm}}.\\
If I reverse the order I get \underline{\hspace{1.5cm}},
so \((f\circ g)(x) \neq (g\circ f)(x)\).''

\vspace{0.15cm}
\small One pair presents --- class signals thumbs.
\end{projcard}

\vspace{0.15cm}
\begin{projcard}{Lesson Wrap}{slideaccent}
Composition order matters in real decisions: discounts, unit conversions, multi-step processes.\\
\(f\circ g \neq g\circ f\) in general --- always verify with a test value.
\end{projcard}
\end{frame}

% ==== Exit Ticket =========================================================
\begin{frame}{Exit Ticket --- Summary Recap}{Algebra 2 \textperiodcentered\ Lesson 5-5 \textperiodcentered\ Period 1}
\phasebadges{1-2}{3}{not graded}
\vspace{0.2cm}
\begin{projcard}{Complete on your exit ticket}{slidegold}
\begin{enumerate}
  \setlength{\itemsep}{0.2cm}
  \item \(f\circ g\) means \underline{\hspace{4cm}} is applied \underline{\hspace{2cm}}.
  \item An example where \(f\circ g \neq g\circ f\): \underline{\hspace{4cm}}.
  \item For a double discount, the better order for the consumer is \underline{\hspace{2cm}} because \underline{\hspace{3cm}}.
  \item One biggest thing I learned today: \underline{\hspace{5cm}}.
\end{enumerate}
\end{projcard}
\end{frame}

\end{document}
